\section{Completion and exact order}\label{sec:completion}

Equations \eqref{eq:sign-bridge} and \eqref{eq:central-identity} imply
\[
 \partial_\xi\Psi(\xi;m,r)<0\qquad(\xi<m<r).
\]
Proposition~\ref{prop:equivalence} therefore proves part~(1) of
Theorem~\ref{thm:main}.

We now prove directly that the classification stops at order four. Let
\(h=211/2000\). Exact rational theta enclosures give
\begin{equation}\label{eq:finite-pf5-witness}
 -1.08\mathbin{\cdot}10^{-7}
 <\det[\Phi((i-j)h)]_{i,j=0}^4
 <-1.05\mathbin{\cdot}10^{-7}.
\end{equation}

\begin{theorem}\label{thm:not-pf5}
The Riemann kernel is not \(\PF_5\).
\end{theorem}

\begin{proof}
The nodes \(0,h,2h,3h,4h\) are strictly increasing. Evenness identifies the
translation matrix with the symmetric Toeplitz matrix determined by the five
kernel values \(\Phi(0),\Phi(h),\ldots,\Phi(4h)\). Each value is enclosed by
exact Taylor bounds for the retained theta modes and the infinite-tail theorem;
a rational determinant parity factorization then gives
\eqref{eq:finite-pf5-witness}. This is a negative translation minor at
distinct nodes.
\end{proof}

Equation~\eqref{eq:finite-pf5-witness} proves part~(2) of
Theorem~\ref{thm:main}. Together with the strict minors through order four,
this completes the exact-order classification.

\begin{remark}[Riemann Hypothesis boundary]
Schoenberg's transform theorem says that the bilateral Laplace transform of a
\(\PF_\infty\) function is the reciprocal of a Laguerre--P\'olya entire function
on its strip of convergence~\cite{schoenberg-ii,belton}. The Fourier transform
of \(\Phi\), up to the standard normalization, is the Riemann \(\Xi\)-function,
which has known real zeros; hence \(\Phi\) is unconditionally not
\(\PF_\infty\). The Schoenberg formulation related to the Riemann Hypothesis
instead asks for the inverse transform of the reciprocal of \(\Xi\) to be a
P\'olya-frequency function~\cite{grochenig}. Thus the finite-order
classification proved here is a different statement.
\end{remark}
